% ======================== ARQUITECTURA PROPUESTA =====================
\section{Arquitectura Propuesta}
\label{sec:arquitectura}

\subsection{Visión General}

La arquitectura es \textbf{por capas}: cada base de datos se alimenta de la anterior y no del archivo de origen. MongoDB recibe el evento limpio; Cassandra construye sus agregados \textit{leyendo MongoDB}; Neo4j construye su grafo \textit{leyendo Cassandra}; y la etapa de indicadores consulta las tres. De este modo MongoDB actúa como única fuente de la verdad y las otras dos bases son vistas derivadas y reconstruibles.

\begin{figure}[H]
    \centering
    \begin{tikzpicture}[
        font=\small,
        box/.style={rectangle, draw=black, thick, minimum width=2.7cm, minimum height=1.15cm, align=center},
        db/.style={cylinder, draw=black, thick, aspect=0.25, shape border rotate=90, minimum width=2.6cm, minimum height=1.9cm, align=center},
        arrow/.style={-{Stealth[scale=1.1]}, thick},
        lbl/.style={font=\scriptsize, align=center}
    ]
    \node[box, fill=gray!12] (src) {\textbf{GH Archive}\\ \texttt{.json.gz} por hora};
    \node[box, fill=green!10, right=1.5cm of src] (etl) {\textbf{Etapas 1--2}\\ Extracción\\ y limpieza};
    \node[db, fill=green!18, right=1.5cm of etl] (mongo) {\textbf{MongoDB}\\ documento\\ original};
    \node[db, fill=blue!14, below=1.5cm of mongo] (cass) {\textbf{Cassandra}\\ agregados\\ y KPI};
    \node[db, fill=violet!16, left=1.5cm of cass] (neo) {\textbf{Neo4j}\\ grafo de\\ relaciones};
    \node[box, fill=orange!14, left=1.5cm of neo] (kpi) {\textbf{Etapa 6}\\ Indicadores};
    \node[box, fill=yellow!18, below=1.9cm of neo, minimum width=5.2cm] (dash) {\textbf{Dashboard Streamlit}\\ consulta las tres bases};

    \draw[arrow] (src) -- node[lbl, above] {descarga} (etl);
    \draw[arrow] (etl) -- node[lbl, above] {etapa 3} (mongo);
    \draw[arrow] (mongo) -- node[lbl, right] {etapa 4} (cass);
    \draw[arrow] (cass) -- node[lbl, above] {etapa 5} (neo);
    \draw[arrow] (neo) -- (kpi);
    \draw[arrow, dashed] (kpi.south) |- node[lbl, pos=0.75, below] {\texttt{kpis\_by\_date}} ($(cass.south)+(0,-0.45)$) -- (cass.south);
    \draw[arrow, gray] (neo.south) -- (dash.north);
    \draw[arrow, gray] (cass.south west) -- (dash.north east);
    \draw[arrow, gray] (mongo.east) -- ++(0.9,0) |- (dash.east);

    \begin{scope}[on background layer]
        \node[draw=black!50, dashed, rounded corners, fit=(etl)(mongo)(cass)(neo)(kpi), inner xsep=0.45cm, inner ysep=0.7cm, label={[font=\scriptsize, anchor=south west]north west:DAG de Apache Airflow (cada hora)}] {};
    \end{scope}
    \end{tikzpicture}
    \caption{Arquitectura por capas. Las flechas negras son escrituras del DAG; las grises, lecturas del \textit{dashboard}. La etapa 6 lee las tres bases y guarda los indicadores en Cassandra.}
    \label{fig:arquitectura}
\end{figure}

\subsection{Justificación Técnica por Base de Datos}
\label{subsec:justificacion}

El enunciado exige justificar por qué cada tipo de dato se almacena en una base determinada. La Tabla~\ref{tab:justificacion} resume la decisión y los párrafos siguientes la sustentan con las características del dato y de las consultas.

\begin{table}[H]
\centering
\caption{Qué almacena cada base de datos y por qué}
\label{tab:justificacion}
\vspace{0.2cm}
\small
\begin{tabular}{@{}p{2.1cm}p{4.3cm}p{8.3cm}@{}}
\toprule
\textbf{Base} & \textbf{Qué almacena} & \textbf{Por qué} \\
\midrule
MongoDB &
Documento JSON de cada evento (\texttt{actor}, \texttt{repo}, \texttt{org}, \texttt{payload}, \texttt{created\_at}). &
El \texttt{payload} tiene una forma distinta para cada uno de los 16 tipos de evento observados. Un esquema flexible lo acepta sin migraciones y permite agregar sobre campos anidados como \texttt{payload.action}. \\
\addlinespace
Cassandra &
Conteos por hora y tipo, actividad diaria por repositorio y por usuario, totales diarios, interacciones usuario--repositorio por hora e indicadores por fecha. &
Son datos \textit{append-only} indexados por tiempo. Con el día como clave de partición, cada gráfico del tablero se resuelve leyendo una sola partición, sin \textit{joins} ni agregaciones en lectura. \\
\addlinespace
Neo4j &
Nodos \texttt{User}, \texttt{Repo} y \texttt{Org}, y relaciones tipadas entre ellos con conteo y fechas de primera y última actividad. &
El grado de un repositorio, su vecindario o los colaboradores a dos saltos se obtienen recorriendo relaciones. En un modelo tabular requerirían auto-\textit{joins} sobre millones de filas. \\
\bottomrule
\end{tabular}
\end{table}

\paragraph{MongoDB: el dato tal como llega.} Un \texttt{PushEvent} trae la referencia de la rama; un \texttt{PullRequestEvent} trae el número, el título y la acción; un \texttt{WatchEvent} casi nada. Forzar esa variedad en columnas fijas obligaría a decenas de columnas mayormente nulas o a una tabla por tipo. El modelo documental conserva cada evento íntegro y es la única capa donde sobrevive el \texttt{payload}; por eso el indicador de tasa de fusión de \textit{pull requests} (KPI~4) solo puede calcularse aquí.

\paragraph{Cassandra: la respuesta precalculada.} El tablero formula siempre las mismas preguntas (eventos por hora de un día, ranking de repositorios de un día, serie diaria). Cassandra invierte el diseño habitual: primero se define la consulta y luego la tabla que la responde. El costo de escritura es bajo y predecible, y la lectura no depende del volumen total almacenado sino del tamaño de una partición.

\paragraph{Neo4j: la relación como dato de primera clase.} Los 2.6 millones de relaciones usuario--repositorio del periodo se almacenan como aristas físicas. Preguntar por los usuarios distintos conectados a un repositorio es contar las aristas entrantes de un nodo, no escanear una tabla de hechos.

\subsection{Componentes y Despliegue}

El sistema se ejecuta en contenedores Docker definidos en un único archivo \repo{docker-compose.yml}{docker-compose.yml}, desplegado en una máquina virtual de Google Cloud (Compute Engine, tipo \texttt{e2-standard-4}: 4~vCPU, 16~GB de RAM y disco de 80~GB, región \texttt{us-central1}). Las tres bases de datos corren en esa máquina y no en servicios administrados; esta decisión evita los límites de los niveles gratuitos (512~MB en MongoDB Atlas, 200,000 nodos en Neo4j AuraDB), que el volumen del proyecto supera ampliamente. La orquestación, en cambio, sí usa un servicio administrado: \textbf{Cloud Composer} ejecuta el DAG cada hora y se conecta a las bases por la red privada de la nube (Sección~\ref{subsec:composer}).

\begin{table}[H]
\centering
\caption{Servicios del sistema: contenedores de \texttt{docker-compose.yml} en la máquina virtual y orquestador administrado}
\label{tab:servicios}
\vspace{0.2cm}
\footnotesize
\begin{tabular}{@{}lllp{5.5cm}@{}}
\toprule
\textbf{Servicio} & \textbf{Imagen} & \textbf{Puerto} & \textbf{Función} \\
\midrule
\texttt{mongo}             & \texttt{mongo:7}           & 27017 & Documentos originales. Límite 1.5~GB, caché WiredTiger 0.5~GB. \\
\texttt{cassandra}         & \texttt{cassandra:4.1}     & 9042  & Agregados e indicadores. Límite 3~GB, \textit{heap} 2~GB. \\
\texttt{neo4j}             & \texttt{neo4j:5}           & 7474, 7687 & Grafo. Límite 2.5~GB, \textit{heap} 1~GB, \textit{page cache} 512~MB. \\
\texttt{airflow-scheduler} & Airflow 2.10.5 (Python 3.11) & --  & Planificador de respaldo en la máquina virtual (\texttt{LocalExecutor}); su DAG permanece en pausa mientras Composer está activo. \\
\texttt{airflow-webserver} & Airflow 2.10.5             & 8080  & Interfaz web del Airflow de respaldo. \\
\texttt{postgres}          & \texttt{postgres:15}       & --    & Metadatos de Airflow. \\
\texttt{jupyter}           & imagen propia              & 8888  & \textit{Notebooks} de CRUD y scripts de carga histórica. \\
\texttt{dashboard}         & Streamlit 1.66             & 8501  & Tablero analítico. \\
\texttt{mongo-express}     & \texttt{mongo-express:1.0.2} & 8081 & Explorador web de MongoDB. \\
\texttt{cassandra-web}     & \texttt{cassandra-web:v1.1.6} & 8083 & Explorador web de Cassandra, en modo de solo lectura. \\
\midrule
Cloud Composer 3           & Airflow 2.11.1 administrado & HTTPS & Planifica y ejecuta el DAG cada hora. Entorno pequeño, un \textit{worker} de 1~vCPU y 4~GB. \\
\bottomrule
\end{tabular}
\end{table}

Los diez contenedores consumen en conjunto alrededor de 7.5~GB de memoria en operación (medición con \texttt{docker stats} en la máquina virtual: Cassandra 2.42~GB, Neo4j 2.15~GB, MongoDB 0.98~GB, los dos procesos de Airflow 1.29~GB y el resto por debajo de 0.4~GB cada uno). Los límites de memoria por servicio están fijados de modo que el conjunto quepa en los 16~GB de la máquina.

\paragraph{Red y exposición.} El único puerto abierto a internet es el 8501 del \textit{dashboard}. Las interfaces de administración escuchan en \texttt{127.0.0.1} y se acceden por túnel SSH: Airflow, Jupyter y un explorador web propio para cada base de datos (Mongo Express para MongoDB, Cassandra Web para Cassandra y Neo4j Browser, incluido en la imagen de Neo4j), de modo que el contenido de las tres puede inspeccionarse directamente, además de a través del \textit{dashboard}. Los puertos de las tres bases de datos se publican además en la dirección interna de la máquina, alcanzable solo desde la red privada del proyecto, que es por donde se conecta Cloud Composer; el cortafuegos de la red no los abre al exterior, lo que se comprobó intentando la conexión desde fuera.

\paragraph{Reproducibilidad.} La carpeta \texttt{gcp/} del repositorio contiene los \textit{scripts} para crear la instancia, instalar Docker, abrir el puerto del \textit{dashboard} y levantar el \textit{stack}. Las cifras de volumen de este informe se midieron en esa máquina virtual.

\paragraph{Configuración.} Todos los parámetros (credenciales, rango de la carga histórica, tope de eventos por hora, retraso respecto a la publicación de la fuente) se leen de variables de entorno en \repo{src/config.py}{src/config.py}, de modo que el DAG, los \textit{notebooks} y el \textit{dashboard} comparten exactamente la misma configuración y las mismas funciones de conexión.
